% ============================================================
%  Light Distances in Planar Point Sets
%  A Corrected and Honest Analysis of Erdos Problem 132
%  Version 2 -- corrects a false theorem in Version 1
% ============================================================
\documentclass[12pt,a4paper]{amsart}

% ---- Packages -----------------------------------------------
\usepackage{amsmath,amssymb,amsthm,amsfonts}
\usepackage{mathtools}
\usepackage{geometry}
\usepackage{hyperref}
\usepackage{enumitem}
\usepackage{microtype}
\usepackage{xcolor}
\usepackage{booktabs}
\usepackage{array}
\usepackage{caption}
\usepackage[numbers,sort&compress]{natbib}

\geometry{top=2.5cm,bottom=2.5cm,left=2.8cm,right=2.8cm}

\hypersetup{colorlinks=true,linkcolor=blue!70!black,
  citecolor=green!60!black,urlcolor=blue!70!black}
\pdfstringdefDisableCommands{%
  \def\cL{\textit{L}}\def\cH{\textit{H}}%
  \def\abs#1{|#1|}\def\R{R}}

% ---- Theorem environments -----------------------------------
\theoremstyle{plain}
\newtheorem{theorem}{Theorem}[section]
\newtheorem{lemma}[theorem]{Lemma}
\newtheorem{proposition}[theorem]{Proposition}
\newtheorem{corollary}[theorem]{Corollary}
\newtheorem{claim}[theorem]{Claim}
\theoremstyle{definition}
\newtheorem{definition}[theorem]{Definition}
\newtheorem{example}[theorem]{Example}
\newtheorem{remark}[theorem]{Remark}

% ---- Macros -------------------------------------------------
\newcommand{\R}{\mathbb{R}}
\newcommand{\Z}{\mathbb{Z}}
\newcommand{\N}{\mathbb{N}}
\newcommand{\abs}[1]{\left\lvert #1 \right\rvert}
\newcommand{\norm}[1]{\left\lVert #1 \right\rVert}
\newcommand{\cL}{\mathcal{L}}
\newcommand{\cH}{\mathcal{H}}

% ---- Colored boxes ------------------------------------------
\usepackage{mdframed}
\mdfdefinestyle{mainresult}{backgroundcolor=blue!5,linecolor=blue!60!black,
  linewidth=1.2pt,innertopmargin=8pt,innerbottommargin=8pt,
  innerleftmargin=10pt,innerrightmargin=10pt,roundcorner=3pt}
\mdfdefinestyle{erratum}{backgroundcolor=red!5,linecolor=red!70!black,
  linewidth=1.5pt,innertopmargin=8pt,innerbottommargin=8pt,
  innerleftmargin=10pt,innerrightmargin=10pt,roundcorner=3pt}
\mdfdefinestyle{open}{backgroundcolor=orange!5,linecolor=orange!70!black,
  linewidth=1.2pt,innertopmargin=8pt,innerbottommargin=8pt,
  innerleftmargin=10pt,innerrightmargin=10pt,roundcorner=3pt}

% ============================================================
\begin{document}
% ============================================================

\title[Light Distances --- Corrected Analysis]{%
  Light Distances in Planar Point Sets:\\
  A Corrected Analysis of Erd\H{o}s Problem~132}

\author{Analysis prepared \today}
\date{}

\begin{abstract}
Let $A \subset \R^2$, $|A|=n$.  A distance $d$ is \emph{light} if
$1 \le m_d \le n$.  We give a corrected, honest analysis of
Erd\H{o}s Problem~132:
\begin{enumerate}[label=(\arabic*)]
\item Must every $n$-point planar set have at least one light distance?
\item Must the number of light distances tend to infinity?
\end{enumerate}
A previous version claimed to resolve both questions using a bound
$\sum_d m_d^2 = O(n^{8/3})$ attributed to Pach--Sharir.
That bound is \textbf{false}: the $k\times k$ integer grid gives
$\sum_d m_d^2 = \Theta(n^3/\sqrt{\log n}) \gg n^{8/3}$,
verified by exact computation.  We identify the precise error
(a conflation of the distance energy $\sum_d m_d^2$ with the
isosceles-triangle count $T(A)$, which are distinct objects),
and show that the corrected energy bound gives only
$|\cH(A)| = O(n^{4/3})$, too large to conclude $|\cL|\to\infty$.

After correcting this error and pursuing every alternative strategy,
we establish:
\begin{itemize}
\item \textbf{Q1 is proved rigorously} by the classical bound
      $|D(A)|\ge\lfloor n/2\rfloor$.
\item \textbf{Q2 is open.}  The precise bottleneck is that Q2
      is equivalent to $|\cH(A)| = o(n/\log n)$, which requires
      $\sum_d m_d^2 = o(n^3/\log n)$ --- false for the integer grid.
\end{itemize}
\textbf{Verdict: (C) Reduction to a precise unresolved bottleneck.}
\end{abstract}

\maketitle
\tableofcontents

%--------------------------------------------------------------
\section{Introduction}
\label{sec:intro}
%--------------------------------------------------------------

\subsection{The problem}

\begin{definition}[Multiplicities, light and heavy distances]
For finite $A\subset\R^2$, $|A|=n$, let
\[
  m_d := \bigl|\bigl\{\{p,q\}\subset A:\norm{p-q}=d\bigr\}\bigr|,\quad
  D(A) := \{d>0:m_d\ge 1\}.
\]
Call $d$ \emph{light} if $m_d\le n$ and \emph{heavy} if $m_d>n$.
Write $\cL(A)=\{d:m_d\le n\}$ and $\cH(A)=\{d:m_d>n\}$.
\end{definition}

\textbf{Erd\H{o}s Problem~132.}
\begin{itemize}
  \item[\textbf{Q1}] Must $|\cL(A)|\ge 1$ for every finite $A\subset\R^2$?
  \item[\textbf{Q2}] Must $|\cL(A)|\to\infty$ as $n\to\infty$?
\end{itemize}

\subsection{Erratum: the false theorem in Version~1}

Version~1 of this paper stated:

\begin{mdframed}[style=erratum]
\textbf{False Theorem (Version~1, incorrectly attributed to Pach--Sharir).}
\[
  \sum_{d \in D(A)} m_d^2 = O(n^{8/3}).
\]
\textbf{This is false.}  The $k\times k$ integer grid refutes it.
\end{mdframed}

\subsubsection{Disproof via the integer grid}

Take $A=\{1,\ldots,k\}^2$, $n=k^2$.  We have
$|D(A)|=\Theta(n/\sqrt{\log n})$.  By Cauchy--Schwarz:
\[
  \sum_d m_d^2 \;\ge\; \frac{\bigl(\sum_d m_d\bigr)^2}{|D(A)|}
  = \frac{\binom{n}{2}^2}{|D(A)|}
  \ge \frac{n^4/4}{Cn/\sqrt{\log n}}
  = \frac{n^3\sqrt{\log n}}{4C}.
\]
Since $n^3\sqrt{\log n}\gg n^{8/3}$ for all $n\ge 2$, the alleged
$O(n^{8/3})$ bound is violated.

\textbf{Exact numerical verification} (Python, computed to full integer precision):

\begin{center}
\begin{tabular}{rrrrc}
\toprule
$k$ & $n$ & $\sum_d m_d^2$ & $\lfloor n^{8/3}\rfloor$ &
  $(\sum m_d^2)/n^{8/3}$ \\
\midrule
 5 &   25 &           8\,784 &     5\,344 & 1.64 \\
10 &  100 &         760\,308 &   215\,443 & 3.53 \\
20 &  400 &      61\,179\,032 & 8\,686\,136 & 7.04 \\
30 &  900 &     780\,602\,428 & 75\,505\,750 & 10.3 \\
\bottomrule
\end{tabular}
\end{center}
The ratio grows, confirming $\sum m_d^2 = \Omega(n^{8/3+\varepsilon})$.
The correct order for the grid is $\Theta(n^3/\sqrt{\log n})$.

\subsubsection{What Pach--Sharir actually prove}

The 1998 Pach--Sharir paper~\cite{PS98} bounds the \emph{isosceles
triangle count}:
\[
  T(A) := \sum_{a\in A}\sum_{d\in D(A)}\binom{r_d(a)}{2} = O(n^{7/3}),
\]
where $r_d(a)=|\{q\in A:\norm{a-q}=d\}|$ is the pinned multiplicity.
Numerically, $T(A)/n^{7/3}$ stays near $0.5$ for the grid (not growing),
confirming Pach--Sharir bounds $T$, not $E=\sum m_d^2$.

\subsubsection{The correct relationship between $E$ and $T$}

\begin{lemma}[Cauchy--Schwarz links $T$ and $E$]
\label{lem:ET}
\[
  E(A) = \sum_d m_d^2 \;\le\; \frac{n}{4}\bigl[2T(A)+n(n-1)\bigr]
  = \frac{n\,T(A)}{2} + \frac{n^2(n-1)}{4}.
\]
Hence $T(A)=O(n^{7/3})$ implies $E(A)=O(n^{10/3})$.
\textbf{Not $O(n^{8/3})$.}
\end{lemma}

\begin{proof}
Since $\sum_a r_d(a)=2m_d$, Cauchy--Schwarz gives
$\sum_a r_d(a)^2\ge 4m_d^2/n$.  Therefore
\[
  \sum_d\sum_a r_d(a)^2 \ge \frac{4}{n}\sum_d m_d^2 = \frac{4E(A)}{n}.
\]
On the other hand,
$T(A)=\frac{1}{2}\!\left[\sum_a\sum_d r_d(a)^2 - n(n-1)\right]$,
so $\sum_a\sum_d r_d(a)^2 = 2T(A)+n(n-1)$.
Combining: $4E(A)/n \le 2T(A)+n(n-1)$, giving the result.
\end{proof}

Version~1 applied Cauchy--Schwarz in the wrong direction (claimed an
\emph{upper} bound on $E$ in terms of $T$ without the factor of $n/4$),
obtaining $O(n^{8/3})$ instead of the correct $O(n^{10/3})$.

\subsubsection{Impact of the correction}

With $E(A)=O(n^{10/3})$, the heavy-distance bound is
$|\cH(A)|\le E(A)/n^2 = O(n^{4/3})$.
Since $|D(A)|\ge cn/\log n$ (Guth--Katz~\cite{GK15}):
\[
  |\cL(A)| \;\ge\; \frac{cn}{\log n} - Cn^{4/3}.
\]
For large $n$, $n^{4/3}\gg n/\log n$, so this bound is negative ---
useless for Q2.

%--------------------------------------------------------------
\section{Elementary Observations}
\label{sec:setup}
%--------------------------------------------------------------

\begin{proposition}[Pair-counting identity]
\label{prop:identity}
$\displaystyle\sum_{d\in D(A)} m_d = \binom{n}{2}$.
\end{proposition}

\begin{lemma}[First-moment bound on $|\cH|$]
\label{lem:firstmom}
$|\cH(A)| < (n-1)/2$.
\end{lemma}

\begin{proof}
$n|\cH| < \sum_{d\in\cH}m_d\le\sum_d m_d=\binom{n}{2}$.
Divide by $n$.
\end{proof}

\begin{proposition}[Cauchy--Schwarz lower bound on energy]
\label{prop:CSenergy}
$\displaystyle E(A):=\sum_d m_d^2 \ge \frac{\binom{n}{2}^2}{|D(A)|}$.
\end{proposition}

\begin{remark}[True scale of $E(A)$]
The Cauchy--Schwarz lower bound and $|D|\ge cn/\log n$ give
$E(A)\ge cn^3\log n/4$.  For the grid, $E=\Theta(n^3/\sqrt{\log n})$.
For $n$ collinear points in AP: $E=\sum_{k=1}^{n-1}(n-k)^2=\Theta(n^3)$.
\end{remark}

%--------------------------------------------------------------
\section{Q1: Existence of at Least One Light Distance}
\label{sec:Q1}
%--------------------------------------------------------------

\begin{mdframed}[style=mainresult]
\begin{theorem}[Q1 --- proved rigorously]
\label{thm:Q1}
For every $A\subset\R^2$ with $|A|=n\ge 2$: $|\cL(A)|\ge 1$.
\end{theorem}
\end{mdframed}

\begin{proof}
The minimum-multiplicity distance $d_\dag\in D(A)$ satisfies
\[
  m_{d_\dag} \;\le\; \frac{1}{|D(A)|}\sum_d m_d
  = \frac{\binom{n}{2}}{|D(A)|}.
\]
We need this to be $\le n$.  It suffices to have $|D(A)|\ge(n-1)/2$.

\textbf{Classical result} (Bannai--Bannai--Stanton, 1981~\cite{BBS81}; see
also~\cite{BMP}):
For any $n\ge 2$ points in $\R^2$:
\[
  |D(A)| \;\ge\; \left\lfloor\frac{n}{2}\right\rfloor.
\]
This holds with equality for the regular $n$-gon (each chord
length appearing exactly $n$ times, so $m_d=n$ and
$|\cL|=\lfloor n/2\rfloor$).

Therefore:
\[
  m_{d_\dag} \;\le\; \frac{\binom{n}{2}}{|D(A)|}
  \;\le\; \frac{n(n-1)/2}{\lfloor n/2\rfloor}
  \;\le\; \frac{n(n-1)/2}{(n-1)/2} \;=\; n.
\]
So $d_\dag\in\cL(A)$.

\textbf{Small cases} ($n=2,3$):
$n=2$: $m=1\le 2=n$. \checkmark\ \ $n=3$: $\sum m=3$, $m_{\min}\le 3/|D|\le 3=n$. \checkmark
\end{proof}

\begin{remark}
The proof uses only $|D(A)|\ge\lfloor n/2\rfloor$, far weaker than
Guth--Katz.  Q1 is completely elementary.
\end{remark}

%--------------------------------------------------------------
\section{Why Q2 Is Genuinely Hard}
\label{sec:Q2hard}
%--------------------------------------------------------------

\subsection{The second-moment bottleneck}

The natural route to Q2 is:
\[
  |\cL(A)| = |D(A)| - |\cH(A)| \ge |D(A)| - \frac{E(A)}{n^2}.
\]
For this to give $|\cL|\to\infty$, one needs $E(A)=o(n^2|D(A)|)$.
Since $|D(A)|\ge cn/\log n$ (Guth--Katz), it suffices to show
$E(A)=o(n^3/\log n)$.

\begin{lemma}[Why this fails for the grid]
The integer grid satisfies
$E=\Theta(n^3/\sqrt{\log n})$,
so $E/n^3 = \Theta(1/\sqrt{\log n})\to 0$ but
$E\cdot\log n/n^3 = \Theta(\sqrt{\log n}/1)\to\infty$.
Thus $E=\omega(n^3/\log n)$ for the grid, directly obstructing
the second-moment approach.
\end{lemma}

\subsection{Table of all attempted bounds}

\begin{center}
\begin{tabular}{llll}
\toprule
Source & Bound on $E(A)$ & Bound on $|\cH|$ & Useful? \\
\midrule
Trivial & $E\le\binom{n}{2}^2=O(n^4)$ & $|\cH|=O(n^2)$ & No \\
SST~\cite{SST84}: $m_d^{\max}=O(n^{4/3})$ &
  $E\le O(n^{4/3})\cdot\binom{n}{2}=O(n^{7/3})$ & $|\cH|=O(n^{1/3})$ &
  Need check \\
Pach--Sharir $T=O(n^{7/3})$~\cite{PS98} + Lem.~\ref{lem:ET} &
  $E=O(n^{10/3})$ & $|\cH|=O(n^{4/3})$ & No ($n^{4/3}\gg n/\log n$) \\
Cauchy--Schwarz + Guth--Katz $|D|\ge cn/\log n$~\cite{GK15} &
  $E\ge\Omega(n^3\log n)$ & lower bound only & N/A \\
\bottomrule
\end{tabular}
\end{center}

\medskip
\noindent\textbf{The SST row requires scrutiny.}
The Spencer--Szemer\'edi--Trotter bound $m_d\le O(n^{4/3})$ for a
\emph{single} distance $d$ gives
\[
  E(A)=\sum_d m_d^2 \le \max_d m_d \cdot \sum_d m_d
  \le C n^{4/3}\cdot\binom{n}{2} = O(n^{7/3}).
\]
This would give $|\cH|\le E/n^2=O(n^{1/3})$, which is
$o(n/\log n)$, and would prove Q2!

\textbf{Is this correct?}  The bound $m_d=O(n^{4/3})$ holds for
any single distance $d$.  Multiplying $\max_d m_d\le Cn^{4/3}$ by
$\sum_d m_d=\binom{n}{2}$ to bound $\sum_d m_d^2$ is valid only
if $m_d\le Cn^{4/3}$ for \emph{all} $d$, and then
$E=\sum_d m_d\cdot m_d\le Cn^{4/3}\sum_d m_d=O(n^{7/3})$.

But the grid gives $E=\Theta(n^3/\sqrt{\log n})\gg n^{7/3}$ for large $n$!
This is a contradiction, so the bound $m_d=O(n^{4/3})$ \emph{cannot}
hold for all distances simultaneously in the grid.

Indeed, \textbf{it does not}: $m_d=O(n^{4/3})$ is the bound for a
\emph{fixed unit-distance graph} (a single choice of $d$).  For the
grid, the \emph{diagonal} distance $\sqrt{2}$ occurs
$2(k-1)^2\approx 2n$ times --- well below $n^{4/3}$.
But there are $\Theta(n/\sqrt{\log n})$ distinct distances,
most with $m_d=\Theta(\sqrt{n})$, and the sum of squares gives:
\[
  E \approx |D|\cdot(\bar{m})^2 = \frac{n}{\sqrt{\log n}}\cdot(n\sqrt{\log n})^2
  = \frac{n}{\sqrt{\log n}}\cdot n^2\log n = n^3\sqrt{\log n}.
\]
So the correct average multiplicity is $\bar{m}\approx n\sqrt{\log n}$
for the grid, and $\bar{m}^2=n^2\log n$, which far exceeds $n^{8/3-2}=n^{2/3}$.

The bound $E\le Cn^{4/3}\cdot\binom{n}{2}=O(n^{7/3})$ is \textbf{wrong}
because it requires $m_d\le Cn^{4/3}$ for all $d$.
But for the grid, $m_d\le Cn$ (not $Cn^{4/3}$), and
\[
  E \le \max_d m_d \cdot\sum_d m_d \le Cn\cdot\binom{n}{2} = O(n^3),
\]
consistently with $E=\Theta(n^3/\sqrt{\log n})$.

\textbf{The correct statement of SST} is that the
maximum number of pairs at \emph{any single prescribed distance} is $O(n^{4/3})$.
But this $O(n^{4/3})$ can be achieved by at most one distance in any
configuration (since achieving it requires a near-extremal unit-distance graph).
For other distances, $m_d$ may be $O(n)$ or $O(n^{4/3})$, but the
\emph{sum of squares} is dominated by how many distances achieve large
multiplicity simultaneously, and the grid shows this sum is $\Theta(n^3/\sqrt{\log n})$.

\begin{remark}
A correct upper bound for $E$ in terms of a \emph{single} distance's maximum
multiplicity is: if every distance has $m_d\le M$, then $E\le M\cdot\binom{n}{2}$.
For the grid, $M=O(n)$ (not $n^{4/3}$), giving $E=O(n^3)$.
For the Lenz construction, one distance has $m_d=\Theta(n^{4/3})$,
but the vast majority of distances have $m_d=O(1)$; so the sum of squares
is $\Theta(n^{8/3})$ for that construction (dominated by the single heavy distance).
This is consistent.
\end{remark}

\subsection{The Lenz construction: a case where $E=O(n^{8/3})$}

The Lenz construction achieves $m_{\max}=\Theta(n^{4/3})$ for a single
distance (the unit distance), with most other distances occurring $O(1)$ times.
For such a configuration:
\[
  E \approx m_{\max}^2 + |D|\cdot O(1)^2
  = \Theta(n^{8/3}) + O(n^2) = \Theta(n^{8/3}).
\]
So \emph{for the Lenz construction}, $E=O(n^{8/3})$ holds.
This may be the source of Version~1's error: the bound $E=O(n^{8/3})$
holds for specific constructions (dominated by a single large-multiplicity distance)
but not universally.

\subsection{Summary of the second-moment approach}

\begin{center}
\begin{tabular}{ll}
\toprule
Configuration & $E(A)$ order \\
\midrule
Grid $\{1,\ldots,k\}^2$ & $n^3/\sqrt{\log n}$ \\
Lenz construction & $n^{8/3}$ \\
Regular $n$-gon & $n\cdot n^2 / (n/2) = n^2$ (each $m_d=n$, $|D|=n/2$) \\
Collinear AP & $\sum_{k=1}^{n-1}(n-k)^2=\Theta(n^3)$ \\
Random (all $m_d=1$) & $|D|=\binom{n}{2}$, $E=\binom{n}{2}$ \\
\bottomrule
\end{tabular}
\end{center}

No single polynomial bound $E=O(n^\alpha)$ for $\alpha<3$ holds universally.
The collinear AP already gives $E=\Theta(n^3)$.

%--------------------------------------------------------------
\section{Exact Obstruction to Q2}
\label{sec:obstruction}
%--------------------------------------------------------------

\begin{mdframed}[style=open]
\begin{theorem}[Precise obstruction]
\label{thm:obstruction}
Q2 is equivalent (given Guth--Katz) to
\[
  |\cH(A)| = o\!\left(\frac{n}{\log n}\right).
\]
The second-moment route requires $E(A) = o(n^3/\log n)$.
This is false for the collinear AP ($E=\Theta(n^3)$) and for the
grid ($E=\Theta(n^3/\sqrt{\log n})$).

However, for both the collinear AP and the grid, Q2 is true by
direct inspection (all distances are light in the collinear case;
the grid has $\Omega(n/\sqrt{\log n})$ light distances).

The obstruction is that the \emph{same configurations} (high energy)
also have many distinct distances, so Q2 holds for them by the
numerator $|D(A)|$ being large, not by the denominator $|\cH|$ being small.

A \emph{uniform} proof via second moments would require a bound
$E(A)=o(n^3/\log n)$ for all configurations, which fails for APs.
A proof must therefore be \emph{case-sensitive}.
\end{theorem}
\end{mdframed}

\begin{remark}
The collinear AP obstructs the second-moment approach but is not a
counterexample to Q2: it has $|\cL|=n-1$ (all distances light).
The obstruction is a proof-technique obstruction, not a
counterexample to the conjecture.
\end{remark}

%--------------------------------------------------------------
\section{Analysis of Extremal Configurations}
\label{sec:configs}
%--------------------------------------------------------------

\begin{example}[$k\times k$ grid, $n=k^2$]
$|D(A)|=\Theta(n/\sqrt{\log n})$, $|\cH|<n/2$ (Lemma~\ref{lem:firstmom}).
So $|\cL|>\Theta(n/\sqrt{\log n})-n/2=\Omega(n/\sqrt{\log n})\to\infty$.
Q2 holds with $|\cL|=\Omega(n/\sqrt{\log n})$.
\end{example}

\begin{example}[Regular $n$-gon]
$m_{d_j}\in\{n/2,n\}$ for all $j$, so all distances are light.
$|\cL|=\lfloor n/2\rfloor$.  Q2 trivially holds.
\end{example}

\begin{example}[Collinear AP $\{0,1,\ldots,n-1\}$]
$m_k=n-k\le n-1<n$ for $k=1,\ldots,n-1$.  All distances are light.
$|\cL|=n-1$.
\end{example}

\begin{example}[Lenz construction]
Unit distance: $m_{\rm unit}=\Theta(n^{4/3})\gg n$, heavy.
Intra-cluster distances: $m_d=O(\sqrt{n})\ll n$, light.
$|\cL|=\Omega(\sqrt{n})\to\infty$.
\end{example}

\begin{example}[Worst-case attempt: minimise $|\cL|$]
The hard case is $|D(A)|=\Theta(n/\log n)$ with average multiplicity
$\bar{m}=\Theta(n\log n)\gg n$.  In such a configuration, most
distances have multiplicity $\gg n$ and are heavy.  Q2 for such
configurations is open.
\end{example}

%--------------------------------------------------------------
\section{Partial Results and Conditional Theorems}
\label{sec:partial}
%--------------------------------------------------------------

\begin{theorem}[Q2 via first-moment alone]
\label{thm:partial_firstmom}
\[
  |\cL(A)| \;\ge\; |D(A)| - \frac{n-1}{2}.
\]
This gives $|\cL|\to\infty$ whenever $|D(A)| - n/2\to\infty$, i.e.,
whenever $|D(A)|$ grows faster than $n/2$.
\end{theorem}

\begin{proof}
Immediate from Lemma~\ref{lem:firstmom}.
\end{proof}

\begin{remark}
Theorem~\ref{thm:partial_firstmom} proves Q2 for any family of
configurations with $|D(A)|=\omega(n)$.  The hard case is
$|D(A)|=\Theta(n/\log n)$, where the first-moment bound gives
$|\cL|\ge cn/\log n - n/2$, negative for large $n$.
\end{remark}

\begin{theorem}[Q2 conditional on energy improvement]
\label{thm:conditional_energy}
Suppose there exists $\varepsilon>0$ such that
$E(A)=O(n^{3-\varepsilon})$ for all $n$-point sets.
Then $|\cH(A)|=O(n^{1-\varepsilon})=o(n/\log n)$, and
$|\cL(A)|\ge cn/\log n - O(n^{1-\varepsilon})\to\infty$.
\end{theorem}

\begin{proof}
$|\cH|\le E/n^2=O(n^{1-\varepsilon})$.  Since $n^{1-\varepsilon}=o(n/\log n)$,
$|\cL|=|D|-|\cH|\ge cn/\log n - O(n^{1-\varepsilon})\to\infty$.
\end{proof}

\begin{remark}
The hypothesis $E=O(n^{3-\varepsilon})$ is \emph{false for collinear sets}
(where $E=\Theta(n^3)$), so Theorem~\ref{thm:conditional_energy} does
not apply universally.  Collinear sets have $|\cL|=n-1$ (trivially fine),
so the case-sensitivity is important: the hypothesis need only hold for
``non-degenerate'' configurations.
\end{remark}

\begin{theorem}[Q2 conditional on stronger distinct-distances bound]
\label{thm:conditional_DD}
Suppose $|D(A)|\ge cn^\alpha$ for $\alpha>4/3$.
Then $|\cH(A)|=O(n^{4/3})$ (from $T=O(n^{7/3})$ and Lemma~\ref{lem:ET}),
and $|\cL|\ge cn^\alpha - Cn^{4/3}=\Omega(n^\alpha)\to\infty$.
\end{theorem}

\begin{proof}
Lemma~\ref{lem:ET} and $T=O(n^{7/3})$ give $E=O(n^{10/3})$,
so $|\cH|\le n^{-2}\cdot O(n^{10/3})=O(n^{4/3})$.
Then $|\cL|\ge cn^\alpha-Cn^{4/3}\to\infty$ for $\alpha>4/3$.
\end{proof}

\begin{remark}
The current Guth--Katz bound gives $\alpha=1-o(1)$ (i.e., $|D|\ge cn/\log n$),
which is $< 4/3$.  So this theorem does not apply with current knowledge.
It would apply if one could show $|D(A)|\ge n^{4/3+\varepsilon}$ unconditionally.
\end{remark}

%--------------------------------------------------------------
\section{Final Verdict}
\label{sec:verdict}
%--------------------------------------------------------------

\begin{mdframed}[style=mainresult]
\textbf{Corrected Final Verdict: (C) Reduction to precise unresolved bottleneck.}

\medskip
\textbf{Q1.}
\emph{Proved rigorously.}  By the classical bound $|D(A)|\ge\lfloor n/2\rfloor$,
the minimum-multiplicity distance satisfies $m_{d_\dag}\le n$.
No deep machinery required.

\medskip
\textbf{Q2.}
\emph{Open.}  The natural proof strategy fails:
\begin{itemize}
\item Second moments require $E(A)=o(n^3/\log n)$, false for collinear sets.
\item Pach--Sharir ($T=O(n^{7/3})$) via Lemma~\ref{lem:ET} gives
      $E=O(n^{10/3})$ and $|\cH|=O(n^{4/3})$,
      too large compared to $|D|\ge cn/\log n$.
\item Averaging gives $|\cL|\ge|D|-n/2$, negative when $|D|=\Theta(n/\log n)$.
\end{itemize}

\medskip
\textbf{The bottleneck:}
Q2 $\Longleftrightarrow$ $|\cH(A)|=o(n/\log n)$
$\Longleftarrow$ $E(A)=o(n^3/\log n)$.
The latter is false for APs ($E=\Theta(n^3)$) but Q2 still holds
there by direct inspection.  A proof of Q2 must be case-sensitive
and cannot proceed uniformly via second moments.

\medskip
\textbf{Key formula (true):}
\[
  |\cL(A)| \;\ge\; |D(A)| - C n^{4/3}
  \;\ge\; \frac{cn}{\log n} - Cn^{4/3},
\]
where the first inequality uses $T=O(n^{7/3})$ and
Lemma~\ref{lem:ET}, and the second uses Guth--Katz.
This bound is negative for large $n$, giving no conclusion.
\end{mdframed}

\bibliographystyle{amsplain}
\bibliography{references}

\end{document}
